% from ega0-19-fr.tex
\oldpage[0\textsubscript{IV-1}]{70}

\noindent régulier (22.5.8)~; le \emph{complété} $\widehat A$ de $A$ est alors
un anneau isomorphe à un anneau de séries formelles
$K[[T_1,\ldots,T_n]]$ (19.6.5). En outre, la structure de $k$-algèbre de $A$,
lorsque $A$ est supposé être \emph{complet} et formellement lisse sur $k$, est
entièrement déterminée par le corps résiduel $K$ de $A$ et la \emph{dimension}
de $A$~; celle-ci peut d'ailleurs être arbitraire à condition de satisfaire à
l'inégalité
\[
  \dim(A)\geq \operatorname{rg}_k(\Gamma_{K/k}),
\]
où $\Gamma_{K/k}$ est le «~module d'imperfection~» de $K$ (22.2.6).

En particulier, pour qu'une \emph{extension} $K$ de $k$ soit formellement
lisse, il faut et il suffit que $K$ soit une extension \emph{séparable} de $k$
(19.6.1).

\medskip\noindent\textbf{III.} --- Soient $A$ un anneau local noethérien,
$\mathfrak J$ un idéal de $A$ distinct de $A$, $A_0=A/\mathfrak J$, $B_0$ un
anneau local noethérien \emph{complet}, $A_0\to B_0$ un homomorphisme local
faisant de $B_0$ une $A_0$-algèbre formellement lisse. Alors il existe un
anneau local noethérien \emph{complet} $B$, un homomorphisme local $A\to B$
faisant de $B$ un $A$-module plat, et un $A_0$-\emph{isomorphisme}
\[
  B\otimes_A A_0\xrightarrow{\sim}B_0
\]
(de sorte que, par I), $B$ est une $A$-algèbre \emph{formellement lisse})~; en
outre $B$ est déterminé par ces propriétés à isomorphisme près (19.7.2). Ce
théorème contient en particulier les \emph{théorèmes de Cohen} sur la structure
des anneaux locaux noethériens \emph{complets} (19.8), qui joueront un rôle
important dans les §§~6 et 7 du chapitre IV.

\medskip\noindent\textbf{IV.} --- L'intérêt de l'étude des anneaux locaux
noethériens formellement lisses sur un autre provient de la caractérisation
«~ponctuelle~» suivante des morphismes lisses~: si $X$ et $Y$ sont des
préschémas localement noethériens, $f:X\to Y$ un morphisme localement de type
fini, alors, pour que $f$ soit lisse, il faut et il suffit que pour tout
$x\in X$, l'anneau $\mathcal O_x$ soit formellement lisse sur
$\mathcal O_{f(x)}$. En particulier, si $Y=\operatorname{Spec}(k)$, où $k$ est
un corps parfait, dire que $f$ est lisse équivaut à dire que $X$ est un
préschéma \emph{régulier}.

\medskip\noindent\textbf{V.} --- Enfin, nous verrons aux §§~20, 21 et 22 que
la notion d'algèbre formellement lisse s'introduit de façon naturelle dans la
théorie des \emph{différentielles de Kähler}, les deux théories s'éclairant
mutuellement.
\end{env}

\begin{env}[19.0.3]
\label{0.19.0.3-fr}
Dans tout ce paragraphe et les suivants, les anneaux et modules topologiques
seront supposés \emph{linéairement topologisés} ($\mathbf{0}_{\mathrm I}$,
7.1.1)~; les anneaux topologiques considérés seront supposés \emph{commutatifs},
sauf mention expresse du contraire. Rappelons que si $A$ et $B$ sont deux
anneaux topologiques, $\rho:A\to B$ un homomorphisme d'anneaux définissant sur
$B$ une structure de $A$-algèbre, on dit que $B$ est une \emph{$A$-algèbre
topologique} si $\rho$ est \emph{continu} pour les topologies envisagées.

Pour abréger, dans un anneau topologique $A$ (resp. un $A$-module topologique
$M$), nous dirons «~\emph{système fondamental d'idéaux} (resp.
\emph{sous-modules}) \emph{ouverts}~» au lieu de «~système fondamental de
voisinages de 0 formé d'idéaux (resp. sous-modules)~».

Étant donnés un anneau topologique $A$ et un $A$-module $M$, les ensembles
$\mathfrak J M$, où $\mathfrak J$ parcourt un système fondamental d'idéaux
ouverts, forment un système fondamental de sous-modules ouverts d'une topologie
sur $M$ faisant de $M$ un $A$-module topologique, et que l'on dit \emph{déduite}
de la topologie de $A$.

Soit $M$ un $A$-module topologique dont la topologie est \emph{moins fine} que
la topologie déduite de celle de $A$~; alors, si $N$ est un sous-module
\emph{ouvert} de $M$, le $A$-module \emph{discret} $M/N$ est \emph{annulé par
un idéal ouvert} de $A$, car par hypothèse il existe un tel idéal $\mathfrak R$
tel que $\mathfrak R.M\subset N$.

